\documentclass[10pt,fleqn,twoside]{book}
\topmargin=3cm

%cambiar fuente defecto
\usepackage{helvet}
%\renewcommand{\normalfont}{\sfdefault}

\usepackage{amsmath, amsfonts, psfrag, fancyhdr, layout, appendix, subfig}
\usepackage{graphicx}

\usepackage[utf8]{inputenc}
\usepackage[spanish]{babel}
\usepackage{makeidx}
%\usepackage{ucs}

%imagenes donde deben ir
\usepackage{float}

\usepackage{color}

\definecolor{lightgray}{rgb}{0.83,0.83,0.83}
\definecolor{darkgray}{gray}{0.40}

%Colores predefinidos para usar con listing
\definecolor{code_fg}{RGB}{50,50,50}
\definecolor{code_bg}{RGB}{240, 255, 225}
\definecolor{code_coment}{RGB}{0, 80, 0} 
\definecolor{code_keyword}{RGB}{0, 50, 100}
\definecolor{code_indetif}{RGB}{0, 128, 255}
\definecolor{code_number}{RGB}{80, 80, 80}
\definecolor{console_bg}{RGB}{0, 79, 117}
\definecolor{console_fg}{RGB}{250, 250, 250}
\definecolor{console_key}{RGB}{0, 79, 157}

%color de los enlaces
\definecolor{link}{RGB}{0, 79, 157}


%This change labels of subfig
\renewcommand{\thesubfigure}{\alph{subfigure}\arabic{subfiggroup}}
\captionsetup[subfigure]{labelformat=simple,labelsep=colon,
                         listofformat=subsimple}
%\captionsetup{lofdepth=2} This is in order to list the subfigures in the LOF
\makeatletter
 \renewcommand{\p@subfigure}{}
  %Esto lo agrego yo para tener subfiguras a1, b1, ... a2, b2, ... 
  %Se reinicia cada vez que una nueva figura es convocada (como es debido).
  \newcounter{subfiggroup}[figure] 
\makeatother

\usepackage{epsfig}

%Esto genera enlaces en el PDF
\usepackage{url}

%cambiar margenes
\def\changemargin#1#2{\list{}{\rightmargin#2\leftmargin#1}\item[]}
\let\endchangemargin=\endlist


% \usepackage{html}
%Esto es para el conversor latex2html


%Definición de margenes
\usepackage[left=4.5cm,top=3cm,right=2cm,bottom=2.5cm]{geometry}
\sloppy
\pagestyle{empty}

% Code for creating empty pages
% No headers on empty pages before new chapter
\makeatletter
\def\cleardoublepage{\clearpage\if@twoside \ifodd\c@page\else
    \hbox{}
    \thispagestyle{plain}
    \newpage
    \if@twocolumn\hbox{}\newpage\fi\fi\fi}
\makeatother \clearpage{\pagestyle{plain}\cleardoublepage}

% Code for creating fully-empty pages
% Fully empty pages before command is called
\makeatletter
\def\clearfullypage{\clearpage\if@twoside \ifodd\c@page\else
    \hbox{}
    \thispagestyle{empty}
    \newpage
    \if@twocolumn\hbox{}\newpage\fi\fi\fi}
\makeatother \clearpage{\pagestyle{empty}\clearfullypage}

% Dutch style of paragraph formatting, i.e. no indents.
\setlength{\parskip}{1.3ex plus 0.2ex minus 0.2ex}
\setlength{\parindent}{0pt}


% Double space for REVISION
%\renewcommand{\baselinestretch}{2.0}

%Print subsubsection numbers and put them in TOC
\setcounter{secnumdepth}{3}
\setcounter{tocdepth}{3}

\makeindex

%enlaces web en pdf remplazado por url
\usepackage{hyperref}

\hypersetup{
    colorlinks=true,
    linkcolor=black,
    urlcolor=link
}


%listings reconoce la sintaxis de varios lenguajes y permite otras configuraciones.
\usepackage{listings}
%listing definicion para formateo de codigo fuente
\lstdefinestyle{Python}
{
    basicstyle=\small\ttfamily\color{code_fg},
    backgroundcolor=\color{code_bg},
    keywordstyle=\color{code_keyword}\tt, % underlined bold black keywords
	identifierstyle=\color{code_indetif}, % color identificadores
	commentstyle=\color{code_coment},                       % comentarios
	stringstyle=\ttfamily,                                 % typewriter type for strings
	showstringspaces=false,
	frame=shadowbox,%\linewidth=3,
	framexleftmargin=8mm,
	rulesepcolor=\color{code_fg},
	backgroundcolor=\color{code_bg},
	language=Python, 
	numbers=left, 
	numberstyle=\color{code_number}\footnotesize\tt,		%estilo numero linea
	tabsize=4, 						                        %tamaño tabulacion 	    
	%showspaces=true,
}

\lstdefinestyle{HTML}
{
    basicstyle=\small\ttfamily\color{code_fg},
    backgroundcolor=\color{code_bg},
    keywordstyle=\color{code_keyword}\tt, % underlined bold black keywords
	identifierstyle=\color{code_indetif}, % color identificadores
}

\lstdefinestyle{consola}
{
    basicstyle=\ttfamily\color{console_fg},
    backgroundcolor=\color{console_bg},
}


%FIN Definicion para Formateo Codigo

\usepackage{lmodern}% http://ctan.org/pkg/lm 

\begin{document}


%Cambiar Cuadros por Tablas y lista de... Debe ir después de \begin{document}
%\renewcommand{\listtablename}{Índice de tablas} 
%\renewcommand{\tablename}{Tabla} 


\pagenumbering{arabic}

%%%%%%%%%%%%%
\frontmatter
%%%%%%%%%%%%%
%\ input{portada} %Se compila aparte y se junta con "gs".
%gs -dNOPAUSE -sDEVICE=pdfwrite -sOUTPUTFILE=tesiscompleta.pdf -dBATCH portada.pdf tesis.pdf
\pagestyle{empty}
\clearfullypage


% Define pagestyle
\pagestyle{fancy}
\fancyhf{}
\renewcommand{\chaptermark}[1]{\markboth{ \emph{#1}}{}}
\fancyhead[LO]{}
\fancyhead[LO]{}
\fancyfoot[LE,RO]{\thepage}

% Redefine plain page style
\fancypagestyle{plain}{
\fancyhf{}
\renewcommand{\headrulewidth}{0pt}
\fancyfoot[LE,RO]{\thepage}
}

% Dutch style of paragraph formatting, i.e. no indents.
\setlength{\parskip}{1.3ex plus 0.2ex minus 0.2ex}


% Remove parskip for toc
\setlength{\parskip}{0ex plus 0.5ex minus 0.2ex}


%Portada
%\documentclass[a4paper,10pt]{report}
%\usepackage{lmodern}% http://ctan.org/pkg/lm 
%\usepackage[spanish]{babel} %espaniol
%\usepackage[latin1]{inputenc} %acentos sin codigo
%\usepackage{graphicx} 
%\topmargin=3cm

%\begin{document}


\begin{titlepage}


\begin{center}
    {\fontsize{35}{35} \selectfont
    Anteproyecto: \\
    Sistema de Tarjetas de fidelizacion de clientes. \\[1.3cm] }
\end{center}


\begin{center}
    \LARGE Universidad Nacional de Salta \\
    \begin{figure}[h]
        \begin{center}
        \includegraphics[scale=0.5]{resourse/logo-UNSa.jpg}
        \end{center}
    \end{figure}

    
    \LARGE Facultad de Ciencias Exactas \\
    Seminario de Sistemas \\ [1.3cm]
\end{center}


\begin{flushright}
    {\Large 
    Alumno: Ricardo Daniel Quiroga \\
    L.U.: 212.625 \\
    Director: Mg. Gustavo Daniel Gil \\
    Numero de Actuacion: 1013/14
    \\[1.0cm]
    \textbf{ Intengrantes Comision Evaluadora } \\
    Lic. Adriana Binda \\
    Lic. Loraine Gimson \\
    Lic. Martin Diaz \\
    }
    
\end{flushright}



\end{titlepage}



%\include{dedicatoria}

\tableofcontents %Indice
\listoffigures %listado de figuras
%\listoftables

%\cleardoublepage
%\ include{resumen} 
%\cleardoublepage
%\ include{resumenen} 
% Adjustments headers
\fancyhead[RO]{\leftmark}

%%%%%%%%%%%%%
\mainmatter
%%%%%%%%%%%%%


% Adjustments headers
\fancyhead[RO]{\leftmark}
\fancyhead[EL]{\emph{Capítulo \thechapter}}
%\setcounter{page}{3}
%\ include{file} busca el archivo file.tex en el directorio actual y se lo tomo como si el contenido de
%file.tex estuviera en vez del comando \include pero insertando una pagina en blanco antes y después del 
%contenido de file.tex por lo que el al escribir file.tex hay que hacer de cuentas que se esta dentro de 
%\begin{document} del documento principal, en este caso este archivo es el principal. \input es un comando 
%similar pero no agrega las paginas en blanco  
%elimine algunos \include por que no aportan nada, el que quedo tiene algunos ejemplos de uso de fancyvrb y listings
%el espacio que queda en \ include es solo para que Texmaker no reconozca el include y lo muestre en la estructura del archivo.

\chapter{Objetivos}

\section{Objectivos Generales}

\section{Objectivos Especificos}


\chapter{Planteamiento del Problema}

 

\chapter{Estudio de Factibilidad}

\section{Factibilidad Tecnica}

\section{Factibilidad Operativa}

\section{Factibilidad Economica}

\section{Factibilidad Legal}



\chapter{Descripci\'on de la Metodologia}



\chapter{Unidades de Operaci\'on}

\chapter{Herramientas a Utilizar}

Una de las herramientas que se utilizar\'a es la Herramienta Case WebRatio que 
es una  herramienta  orientada  justamente  a  WebML (la  metodolog\'{\i}a  que  
estudiar\'e  y aplicar\'e durante el desarrollo del presente seminario) la 
misma se la puede descargar desde el sitio 
\url{http://www.webratio.com/portal/content/es/descargas}. 

Esta herramienta permite desarrollar todos los modelos que la metodolog\'{\i}a 
exige, adem\'as brinda la opci\'on de generar la aplicaci\'on que es algo muy 
interesante ya que se puede observar si en realidad los modelos reflejan lo 
que se pretende tanto en la parte de interfaz como en la funcional. 


Esta  herramienta  CASE para poder funcionar requiere:  

\begin{itemize}
    \item Sistema Operativo: Windows 7, Windows XP, Mac OS X, Linux
    \item Memoria RAM de 512 Mb o superior
    \item Espacio Minimo en el Disco 1 Gb o superior
    \item Resolucion de Pantalla 1024x768 px 
    \item Java Development Kit JDK JRockit 1.6.0\_20 o superior
    \item Apache Tomcat
    \item (Opcional) Plugin Apache Tomcat para WebRatio basado en Eclipse 
\end{itemize}


Para mayor informacion acerca de la instalacion de WebRatio puede consultar 
cite{Instalacion WebRatio}.


Es muy conveniente la utilizaci\'on de esta herramienta por el simple hecho de 
que fue  desarrollada  por  el  grupo  precursor  de  esta  metodolog\'{\i}a,  
adem\'as  de  ser  una herramienta bastante intuitiva y de f\'acil 
utilizaci\'on. 


Tambi\'en se requiere utilizar el servidor Web  Apache Tomcat ya que WebRatio 
genera paginas JSP y son interpretadas por este servidor. Se puede instalar
junto con WebRatio. 

\chapter{Indice Tentativo de la Investigacion}

\section{Calendario de Tareas}

\section{Diagrama de Gant}





% \begin{thebibliography}{X}
    \bibitem{Factibilidad} \textsc{Wikipedia} \textit{Definicion Factibilidad}
     http://es.wikipedia.org/wiki/Factibilidad

%    \bibitem{Apache} \textsc{Wikipedia} \textit{Servidor HTTP Apache} http://es.wikipedia.org/wiki/Servidor\_HTTP\_Apache 

     \bibitem{Instalacion WebRatio} \textsc{Agencia de Informatica y 
     Comunicaciones de la Comunidad de Madrid} \textit{WebRatio Guia de 
     Preparaci\'on de Entorno de Desarrollo}  
     http://www.madrid.org/arquitecturasw/images/documentacion/webratio/actual/
     Entorno/WR\_MUS\_Preparacion\_Entorno\_Desarrollo.pdf
    

    %\bibitem{} \textsc{} \textit{}
\end{thebibliography}


%\part*{\addcontentsline{toc}{part}{Apéndices}Apéndices}



% Adjustments headers
%\fancyhead[RO]{\leftmark}
%\fancyhead[EL]{\emph{Apéndice \thechapter}}

%%%%%%%%%%%%%
%\appendix
%%%%%%%%%%%%%%
%\ include{software}
%\ include{publicaciones}


%%%%%%%%%%%%%
\backmatter
%%%%%%%%%%%%%
% Adjustments headers
\fancyhead[RO]{\leftmark}
\fancyhead[EL]{}
\addcontentsline{toc}{chapter}{Bibliografía}
%\bibliographystyle{unsrt}
%\bibliographystyle{plain}
%\bibliography{biblio.bib}

\include{bibligrafia}

\end{document}

%%% Local Variables: 
%%% mode: latex
%%% TeX-master: "tesis"
%%% End: 
